\documentclass[12pt, letterpaper]{article}

\usepackage[utf8]{inputenc}
\usepackage[T1]{fontenc}
\usepackage[spanish]{babel}
\usepackage{geometry}
\usepackage{amsmath, amssymb, amsthm}
\usepackage{booktabs}
\usepackage{array}
\usepackage{microtype}
\usepackage{setspace}
\usepackage{parskip}
\usepackage{titlesec}
\usepackage{fancyhdr}
\usepackage{enumitem}
\usepackage{bm}
\usepackage{mathtools}
\usepackage{tcolorbox}
\usepackage{hyperref}

\tcbuselibrary{skins, breakable}

\geometry{top=2.8cm, bottom=3.0cm, left=3.2cm, right=3.2cm, headheight=14pt}
\setstretch{1.20}

\pagestyle{fancy}
\fancyhf{}
\fancyhead[C]{\small\textsc{Econometría --- Gujarati --- Ejercicio 13.3}}
\fancyfoot[C]{\small\thepage}
\renewcommand{\headrulewidth}{0.4pt}
\renewcommand{\footrulewidth}{0pt}

\titleformat{\section}
  {\large\bfseries\scshape}{\thesection.}{0.6em}{}
  [\vspace{-4pt}\rule{\linewidth}{0.4pt}]
\titleformat{\subsection}{\normalsize\bfseries}{\thesubsection.}{0.5em}{}
\titleformat{\subsubsection}{\normalsize\itshape}{}{0em}{}

\newtcolorbox{clave}{
  colback=white, colframe=black,
  boxrule=0.5pt, arc=3pt,
  left=8pt, right=8pt, top=6pt, bottom=6pt,
  breakable
}

\newcommand{\Var}{\operatorname{Var}}

\begin{document}

\begin{center}
  {\LARGE\bfseries Ejercicio 13.3 --- Omisión del Intercepto}\\[6pt]
  {\LARGE\bfseries cuando el Modelo Verdadero lo Incluye}\\[12pt]
  {\large\itshape Consecuencias del error de especificación por omisión}\\
  {\large\itshape de una variable relevante: el término constante}\\[14pt]
  \rule{0.55\linewidth}{0.6pt}\\[8pt]
  {\normalsize Ejercicio~13.3 --- \textit{Econometría}, Gujarati \& Porter, 5.a ed.\
  por Emanuel Quintana Silva}\\[4pt]
  {\small\textsc{Verificación empírica con datos simulados, $n = 50$, semilla 12345}}
\end{center}

\vspace{10pt}

%------------------------------------------------------------------------------
\section{Planteamiento del problema}
%------------------------------------------------------------------------------

Se continúa con el ejercicio 13.2, pero ahora el modelo \textbf{(2) es el
verdadero}:

\begin{align}
  \textbf{Modelo verdadero:} \quad
  Y_i &= \alpha_0 + \alpha_1 X_i + v_i \label{eq:verdadero}\\[6pt]
  \textbf{Modelo mal especificado:} \quad
  Y_i &= \beta_1 X_i + u_i \label{eq:origen}
\end{align}

El modelo~\eqref{eq:origen} obliga a la línea de regresión a pasar por el
origen, omitiendo el término constante $\alpha_0 \neq 0$ que sí existe en
la realidad.

En la taxonomía de Gujarati (sección~13.2), este error corresponde a la
\textbf{omisión de una variable relevante}, donde la ``variable'' omitida
es el término constante $\alpha_0$. Es el caso opuesto al ejercicio anterior:

\begin{center}
\renewcommand{\arraystretch}{1.3}
\small
\begin{tabular}{lll}
\toprule
\textbf{Tipo de error} & \textbf{Consecuencia principal} & \textbf{Ejercicio} \\
\midrule
Inclusión de irrelevante & Insesgadez + ineficiencia & 13.2 \\
\textbf{Omisión de relevante} & \textbf{Sesgo + inconsistencia} & \textbf{13.3} \\
\bottomrule
\end{tabular}
\end{center}

%------------------------------------------------------------------------------
\section{Consecuencias teóricas}
%------------------------------------------------------------------------------

\subsection{Sesgo e inconsistencia del estimador de la pendiente}

El estimador MCO de $\beta_1$ en el modelo~\eqref{eq:origen}, forzado a
pasar por el origen, es:

\begin{equation}
  \hat{\beta}_1 = \frac{\sum X_i Y_i}{\sum X_i^2}
\end{equation}

Sustituyendo el modelo verdadero $Y_i = \alpha_0 + \alpha_1 X_i + v_i$:

\begin{align}
  E(\hat{\beta}_1)
  &= \frac{\sum X_i\,E(Y_i)}{\sum X_i^2}
   = \frac{\sum X_i(\alpha_0 + \alpha_1 X_i)}{\sum X_i^2} \notag \\[4pt]
  &= \alpha_1 + \alpha_0\left(\frac{\sum X_i}{\sum X_i^2}\right)
\end{align}

El sesgo resultante es:

\begin{equation}
  \boxed{
    \operatorname{Sesgo}(\hat{\beta}_1)
    = \alpha_0 \cdot \frac{\sum X_i}{\sum X_i^2}
  }
\end{equation}

El estimador $\hat{\beta}_1$ solo sería insesgado si:

\begin{itemize}[leftmargin=1.8em]
  \item $\alpha_0 = 0$ \quad (contradice el supuesto del modelo verdadero), o
  \item $\sum X_i = 0$ \quad (raramente plausible con datos económicos, donde $X_i > 0$).
\end{itemize}

Dado que ninguna de estas condiciones se cumple en general, el sesgo
\textbf{no desaparece al aumentar $n$}: el estimador también es
\textbf{inconsistente}.

\begin{clave}
  A diferencia del Ej.~13.2, donde $E(\hat{\alpha}_1) = \beta_1$ se
  preservaba, aquí $E(\hat{\beta}_1) \neq \alpha_1$ de forma sistemática.
  La dirección del sesgo depende del signo de $\alpha_0$: si $\alpha_0 > 0$
  y $\sum X_i > 0$ (el caso típico), el sesgo es \textbf{positivo} y se
  sobreestima la pendiente.
\end{clave}

\subsection{Estimación incorrecta de la varianza del error}

El estimador de $\sigma^2$ en el modelo~\eqref{eq:origen} usa $n - 1$
grados de libertad:

\begin{equation}
  \hat{\sigma}^2_{(1)} = \frac{\sum \hat{u}_i^2}{n - 1}
\end{equation}

Al forzar la línea a pasar por el origen, los residuos $\hat{u}_i$ deben
absorber el efecto del intercepto omitido $\alpha_0$, produciendo residuos
sistemáticamente más grandes que en el modelo correcto. Esto conduce a una
\textbf{sobreestimación} de $\sigma^2$ y, en consecuencia, a varianzas
estimadas incorrectas para todos los coeficientes.

\subsection{Inferencia estadística inválida}

Dado que tanto los coeficientes como sus varianzas estimadas están sesgados:

\begin{itemize}[leftmargin=1.8em]
  \item Las pruebas $t$ y $F$ \textbf{dejan de ser válidas} y pueden llevar
    a conclusiones erróneas sobre la significancia de $X$.
  \item Los intervalos de confianza \textbf{no tienen la cobertura nominal}:
    no contendrán el verdadero $\alpha_1$ con la probabilidad deseada $(1-\alpha)$.
  \item Los errores de Tipo~I y Tipo~II se cometen de forma sistemática,
    no aleatoria.
\end{itemize}

\subsection{$R^2$ potencialmente negativo}

En regresión por el origen, los residuos no necesariamente suman cero
($\sum \hat{u}_i \neq 0$), por lo que la descomposición estándar de la
varianza total no se sostiene:

\begin{equation}
  \underbrace{\sum(Y_i - \bar{Y})^2}_{SCT}
  \;\neq\;
  \underbrace{SCM}_{\text{explicada}}
  +
  \underbrace{SCR}_{\text{residual}}
\end{equation}

El $R^2$ convencional puede ser \textbf{negativo}, indicando que el modelo
ajustado predice peor que el simple promedio $\bar{Y}$. Stata reporta
$R^2$ no centrado (\texttt{r2\_u}) para \texttt{noconstant}, que no es
comparable con el $R^2$ centrado del modelo con intercepto.

\subsection{Predicciones sesgadas}

Los pronósticos $\hat{Y}_i = \hat{\beta}_1 X_i$ estarán sistemáticamente
desplazados respecto a los valores reales. El sesgo es especialmente
pronunciado para valores de $X$ cercanos a cero, donde el intercepto
omitido tiene mayor peso relativo.

%------------------------------------------------------------------------------
\section{Resumen de consecuencias}
%------------------------------------------------------------------------------

\begin{center}
\renewcommand{\arraystretch}{1.4}
\small
\begin{tabular}{p{4.2cm}p{4.0cm}p{4.2cm}}
\toprule
\textbf{Propiedad}
  & \textbf{Ej.~13.2} \newline (intercepto innecesario)
  & \textbf{Ej.~13.3} \newline (intercepto omitido) \\
\midrule
Insesgadez          & \checkmark\ Preservada        & $\times$\ Perdida \\
Consistencia        & \checkmark\ Preservada        & $\times$\ Perdida \\
Eficiencia          & $\sim$\ Reducida              & --- No aplica \\
Validez $t$ y $F$   & \checkmark\ Válidas           & $\times$\ Inválidas \\
Intervalos de conf. & $\sim$\ Más amplios           & $\times$\ Cobertura incorrecta \\
$R^2$ comparable    & $\sim$\ No comparable         & $\times$\ Puede ser negativo \\
Predicciones        & \checkmark\ Insesgadas        & $\times$\ Sesgadas sistemáticamente \\
\bottomrule
\end{tabular}
\end{center}

%------------------------------------------------------------------------------
\section{Verificación empírica con datos simulados}
%------------------------------------------------------------------------------

Se generan datos del proceso verdadero $Y_i = 3 + 2X_i + v_i$
($\alpha_0 = 3$, $\alpha_1 = 2$, $v_i \sim N(0,1)$, $n = 50$) y se
comparan ambas especificaciones.

\subsection{Modelo correcto --- con intercepto}

\begin{align}
  \widehat{Y}_i &= 2.8711 + 2.0160\,X_i
\end{align}

\begin{center}
\renewcommand{\arraystretch}{1.3}
\small
\begin{tabular}{lrrrr}
\toprule
\textbf{Variable} & \textbf{Coeficiente} & \textbf{Error est.} & \textbf{Razón $t$} & \textbf{$P(>|t|)$} \\
\midrule
$X_i$      &  2.0160 & 0.0491 &  41.02 & 0.000 \\
Constante  &  2.8711 & 0.2928 &   9.81 & 0.000 \\
\midrule
\multicolumn{5}{l}{$R^2 = 0.9723$\quad $\bar{R}^2 = 0.9717$\quad $F(1,48) = 1682.58$\quad $n = 50$} \\
\bottomrule
\end{tabular}
\end{center}

\subsection{Modelo mal especificado --- regresión por el origen}

\begin{align}
  \widehat{Y}_i &= 2.4423\,X_i
\end{align}

\begin{center}
\renewcommand{\arraystretch}{1.3}
\small
\begin{tabular}{lrrrr}
\toprule
\textbf{Variable} & \textbf{Coeficiente} & \textbf{Error est.} & \textbf{Razón $t$} & \textbf{$P(>|t|)$} \\
\midrule
$X_i$  &  2.4423 & 0.0393 & 62.16 & 0.000 \\
\midrule
\multicolumn{5}{l}{$R^2_u = 0.9875$ (no centrado)\quad $F(1,49) = 3863.98$\quad $n = 50$} \\
\bottomrule
\end{tabular}
\end{center}

\begin{clave}
  El $R^2 = 0.9875$ del modelo por el origen es \textbf{no centrado} y no
  es comparable con el $R^2 = 0.9723$ del modelo con intercepto. La
  apariencia de ``mejor ajuste'' en el modelo restringido es un artefacto
  numérico, no evidencia de superioridad.
\end{clave}

%------------------------------------------------------------------------------
\section{Verificación de la fórmula del sesgo}
%------------------------------------------------------------------------------

Con los resultados de Stata se verifica algebraicamente la fórmula del sesgo:

\begin{center}
\renewcommand{\arraystretch}{1.3}
\small
\begin{tabular}{lr}
\toprule
\textbf{Cantidad} & \textbf{Valor} \\
\midrule
$\alpha_0$ verdadero                              &  3.0000 \\
$\sum X_i$                                        & 263.55 \\
$\sum X_i^2$                                      & 1774.68 \\
$\sum X_i / \sum X_i^2$                           &  0.148503 \\
Sesgo teórico $= \alpha_0 \cdot \frac{\sum X_i}{\sum X_i^2}$ & 0.4455 \\
\midrule
$\hat{\alpha}_1$ (modelo correcto)                &  2.0160 \\
$\hat{\beta}_1$  (modelo por el origen)           &  2.4423 \\
Sesgo empírico $= \hat{\beta}_1 - \hat{\alpha}_1$ &  0.4264 \\
Diferencia (empírico $-$ teórico)                 & $-$0.0191 \\
\bottomrule
\end{tabular}
\end{center}

La pequeña diferencia de $-0.0191$ entre el sesgo empírico ($0.4264$) y el
teórico ($0.4455$) se debe a que $\hat{\alpha}_1 \neq \alpha_1 = 2$ exactamente
por variación muestral. La fórmula se verifica con alta precisión.

\begin{clave}
  $\hat{\beta}_1 = 2.4423 \neq 2.0160 \approx \hat{\alpha}_1$. La
  sobreestimación de la pendiente en el modelo restringido
  ($\Delta = 0.4264$) coincide estrechamente con el sesgo teórico predicho
  ($0.4455 = 3 \times 263.55/1774.68$). El intercepto estimado en el modelo
  correcto ($\hat{\alpha}_0 = 2.871$, $t = 9.81$, $p < 0.001$) confirma que
  el término constante era \textbf{necesario} y estadísticamente significativo.
\end{clave}

%------------------------------------------------------------------------------
\section{Tabla comparativa de estimadores}
%------------------------------------------------------------------------------

\begin{center}
\renewcommand{\arraystretch}{1.3}
\small
\begin{tabular}{lcc}
\toprule
                              & \textbf{Modelo correcto} & \textbf{Modelo origen} \\
\textbf{Variable}             & $(\alpha_0, \alpha_1 X)$   & $(\beta_1 X)$         \\
\midrule
$X_i$ (pendiente)             & 2.0160  & 2.4423  \\
Error estándar                & 0.0491  & 0.0393  \\
Constante                     & 2.8711  & ---     \\
Error estándar                & 0.2928  & ---     \\
\midrule
$n$                           & 50      & 50      \\
$R^2$                         & 0.9723  & 0.9875* \\
Root MSE                      & 0.9650  & 1.6552  \\
\midrule
Sesgo de $\hat{\beta}_1$      & 0       & $+0.4264$ \\
\bottomrule
\multicolumn{3}{l}{\footnotesize * $R^2$ no centrado; no comparable con el $R^2$ centrado del modelo con intercepto.}
\end{tabular}
\end{center}

Nótese que el Root MSE del modelo por el origen ($1.6552$) es
\textbf{considerablemente mayor} que el del modelo correcto ($0.9650$),
evidenciando la sobreestimación de $\sigma^2$ predicha teóricamente:
los residuos deben absorber el efecto del intercepto omitido.

%------------------------------------------------------------------------------
\section{Síntesis: comparación de los dos errores de especificación}
%------------------------------------------------------------------------------

\begin{center}
\renewcommand{\arraystretch}{1.4}
\small
\begin{tabular}{p{3.8cm}p{3.8cm}p{3.8cm}}
\toprule
\textbf{Criterio}
  & \textbf{Ej.~13.2} \newline (intercepto innecesario)
  & \textbf{Ej.~13.3} \newline (intercepto omitido) \\
\midrule
Error cometido        & Incluir $\alpha_0$ cuando $\alpha_0 = 0$ & Omitir $\alpha_0$ cuando $\alpha_0 \neq 0$ \\
Tipo en taxonomía     & Variable irrelevante incluida            & Variable relevante omitida \\
$E(\hat{\beta}_1)$    & $= \beta_1$ (insesgado)                  & $\neq \alpha_1$ (sesgado) \\
$\hat{\sigma}^2$      & Levemente inflada                        & Fuertemente sobreestimada \\
Pruebas $t$, $F$      & Válidas                                  & Inválidas \\
Costo principal       & Menor eficiencia                         & Sesgo + inferencia inválida \\
Gravedad              & Leve                                     & \textbf{Grave} \\
\bottomrule
\end{tabular}
\end{center}

%------------------------------------------------------------------------------
\section{Conclusión}
%------------------------------------------------------------------------------

\begin{clave}
  Omitir un intercepto relevante es un error de especificación
  \textbf{mucho más grave} que incluir uno irrelevante. Mientras que el
  segundo (Ej.~13.2) solo genera \textit{ineficiencia} --- estimadores
  insesgados pero con mayor varianza ---, el primero invalida las
  propiedades fundamentales del estimador MCO: se pierde la
  \textbf{insesgadez} y la \textbf{consistencia}, la inferencia estadística
  deja de ser válida y las predicciones estarán sistemáticamente sesgadas.

  Los datos de Stata confirman la teoría con precisión:
  $\hat{\beta}_1 = 2.4423$ sobreestima la pendiente verdadera en
  $+0.4264$, cifra muy cercana al sesgo teórico de $0.4455$. El Root MSE
  se incrementa de $0.9650$ a $1.6552$ al restringir el modelo, evidencia
  directa de la sobreestimación de $\sigma^2$.

  La regla práctica de Gujarati cobra aquí su pleno sentido: \textit{ante la
  duda sobre si incluir o no un intercepto, es preferible incluirlo y
  asumir la pequeña pérdida de eficiencia}. Omitirlo cuando realmente
  existe en el proceso generador de datos produce un error sistemático que
  no se corrige con muestras más grandes.
\end{clave}

\vspace{14pt}
\noindent\rule{\linewidth}{0.4pt}

\noindent{\small\textit{Nota metodológica.}
La asimetría entre los dos errores de especificación tiene implicaciones
prácticas directas. En la estrategia de modelación econométrica, la
\textit{regla general} es partir del modelo más general (con intercepto)
y reducirlo solo si la teoría económica y las pruebas estadísticas lo
justifican conjuntamente. Restringir el modelo a priori sin fundamento
teórico sólido expone al investigador al sesgo documentado en este
ejercicio. Los únicos casos donde la regresión por el origen es
apropiada son aquellos en que la teoría garantiza inequívocamente que
$Y = 0$ cuando $X = 0$ (ej.\ ingresos nulos implican consumo nulo en
ciertos modelos), y aun en estos casos la restricción debería
contrastarse formalmente.}

\vspace{14pt}

\begin{center}
\small
\textsc{Referencias}\\[6pt]
\end{center}

\noindent Gujarati, D.~N., \& Porter, D.~C. (2010). \textit{Econometría}
(5.ª~ed.). McGraw-Hill. Capítulo~13, secciones~13.2 y~13.3.

\noindent Greene, W.~H. (2018). \textit{Econometric Analysis} (8.ª~ed.).
Pearson. Capítulo~5.

\noindent Wooldridge, J.~M. (2016). \textit{Introductory Econometrics}
(6.ª~ed.). Cengage. Capítulo~3.

\end{document}
